\section{问题分析与总结}
第一组围绕人形抱物和地图避障跟随建立了可分步调试的实现。动作层通过固定抓取末帧上肢、复用原步态下肢实现抱物前进与转向；视觉层从开环次数控制发展到道路几何居中及固定标记航路，减轻直行偏航的累计影响。方式二将AR相对观测变换到地图，利用历史路径回退生成目标，再交由导航栈规划，并以任务、定位和扫描等有效性条件约束速度输出。

实验过程表明，影响系统的因素不仅是控制算法，还包括动作版本、抱持负载、相机姿态、USB供电与连接、深度扫描覆盖、多机TF与网络。仅根据模型镜头高度推断障碍可见、仅根据动作次数认定转角或仅根据单帧识别认定完整跟随成功，都会使报告结论超出证据。

新增比赛视频支持了人形持物、通道运行及标记调试的部分现场过程；由于车辆身份未知且终点放置与末车停稳未形成连续证据，仍不填入完整成功率。当前主要待验收项为最新标记路线的完整实机验收、抓取/掉落/放置反馈的一致接入、严格的丢目标原地恢复、任务心跳适配、矮障碍检测复核及全队同场结果统计。后续应以日志和连续视频建立可复核记录，再填写性能结果。其余小组应在各自章节补足实际实现与方式独立性说明，共同完成终稿。
